\chapter{Erythema Multiforme}
\index{Erythema Multiforme}
\label{sec:Erythema Multiforme}

\section{Overview}
Erythema multiforme (EM) is an acute, usually self-limited, hypersensitivity reaction of the skin and mucous membranes, most commonly triggered by infections, particularly herpes simplex virus. It is characterized by target lesions with a distinctive three-zone appearance. EM is classified as EM minor (typical target lesions, minimal mucosal involvement) and EM major (extensive mucosal involvement).
\section{Classification}

\begin{description}[font=\normalfont\bfseries,leftmargin=3em,labelwidth=2.5em]
  \item[Cause] Hypersensitivity / Immune
  \item[Organ System] Skin
  \item[Pathophysiology] Immune-mediated reaction with cytotoxic T-cell response directed against keratinocytes expressing viral antigens
  \item[Duration] Acute (1--6 weeks)
  \item[Onset] Sudden
  \item[Mode of Transmission] Non-communicable
  \item[Clinical Course] Acute, self-limited
  \item[Severity] Mild to Moderate
  \item[Extent of Disease] Widespread
  \item[Age Group] Young adults (20--40 years)
  \item[Epidemiology] Incidence 1--5 per million per year
  \item[ICD-11 Code] EB01.0
\end{description}

\section{Causes}
Herpes simplex virus (HSV-1 and HSV-2) infection is the most common trigger, accounting for 50--70\% of cases, particularly recurrent EM. Other triggers include Mycoplasma pneumoniae infection, other viral infections (EBV, VZV, CMV, hepatitis viruses), medications (particularly NSAIDs, sulfonamides, anticonvulsants, antibiotics), and rarely vaccines. The reaction involves HSV DNA fragments in keratinocytes triggering a T-cell-mediated immune response.

\section{Risk Factors}

\begin{itemize}[noitemsep]
  \item Recurrent herpes simplex virus infection
  \item Recent viral or bacterial infection (especially Mycoplasma pneumoniae)
  \item HLA subtype DQ3 (genetic susceptibility)
  \item Medication use
  \item Immunosuppression
\end{itemize}

\section{Signs and Symptoms}

\subsection{Early symptoms}
\begin{itemize}[noitemsep]
  \item Early manifestations of erythema multiforme may be subtle and nonspecific
\end{itemize}

\subsection{Common symptoms}
\begin{itemize}[noitemsep]
  \item \subsection{Early symptoms}
  \item \subsection{Common symptoms}
  \item \textbf{Common symptoms:}
  \item Sudden onset of red papules that evolve into target lesions with three zones
  \item Classic target lesions have a central blister or crust (dark), middle pale edematous ring, and outer red rim
  \item Lesions are acrally distributed: hands, feet, extensor surfaces of arms and legs
  \item Mucosal involvement in EM major: oral, ocular, genital erosions
  \item Lesions may be pruritic or burning
  \item Systemic symptoms: low-grade fever, malaise, headache
  \item Pain or discomfort in the affected area
  \item Difficulty performing daily activities
\end{itemize}

\subsection{Less common symptoms}
\begin{itemize}[noitemsep]
  \item Atypical or less common presentations of erythema multiforme may occur
\end{itemize}

\subsection{Emergency warning signs}
\begin{itemize}[noitemsep]
  \item Severe or worsening symptoms of erythema multiforme require immediate medical evaluation
\end{itemize}
\section{Progression and Stages}
Lesions develop rapidly over 1--3 days and evolve from macules to papules to target lesions. New lesions may continue to appear for 3--5 days. The eruption typically lasts 1--6 weeks and resolves spontaneously without scarring. Recurrent EM, usually triggered by HSV, tends to follow a milder, shorter course with each episode.

\section{Complications}

\begin{itemize}[noitemsep]
  \item Secondary bacterial infection of eroded skin
  \item Post-inflammatory hyperpigmentation or hypopigmentation
  \item Ocular complications (conjunctivitis, keratitis, symblepharon) in EM major
  \item Esophageal or genital strictures from mucosal involvement
  \item Recurrent episodes (HSV-associated)
  \item Reduced quality of life
  \item Emotional and psychological effects
\end{itemize}

\section{Diagnosis}

\begin{itemize}[noitemsep]
  \item Clinical diagnosis based on characteristic target lesions and acral distribution
  \item Skin biopsy of early lesion shows interface dermatitis with necrotic keratinocytes
  \item Direct immunofluorescence to differentiate from autoimmune blistering diseases
  \item HSV PCR or serology to identify viral trigger
  \item Mycoplasma serology in children
\end{itemize}

\section{Prevention}

\begin{itemize}[noitemsep]
  \item Suppressive antiviral therapy (acyclovir, valacyclovir) for HSV-associated recurrent EM
  \item Avoid known medication triggers
  \item Prompt treatment of HSV outbreaks
  \item Go for regular check-ups so any problems are caught early
\end{itemize}

\section{Treatment Options}

\begin{itemize}[noitemsep]
  \item Supportive care and symptomatic treatment for mild EM minor
  \item Topical corticosteroids for pruritus and inflammation
  \item Oral antihistamines for pruritus
  \item Systemic corticosteroids for severe or extensive disease (controversial)
  \item Antiviral therapy (acyclovir) if active HSV infection present
  \item Oral antibiotics if Mycoplasma is the trigger
  \item Wound care for eroded skin and mucosal lesions
  \item Lifestyle changes and self-care
\end{itemize}

\section{Living with It}

\begin{itemize}[noitemsep]
  \item Follow your treatment plan as prescribed
  \item Keep up with regular appointments with your healthcare provider
  \item Adjust your daily habits as needed
  \item Reach out to family, friends, or support groups for help
  \item Keep learning about your condition and how to manage it
\end{itemize}

\section{Outlook}
EM minor is a self-limited condition with an excellent prognosis, typically resolving within 2--6 weeks without scarring. EM major may require hospitalization for supportive care but also has a good prognosis in most cases. Recurrent HSV-associated EM can be effectively managed with suppressive antiviral therapy. Mortality is rare and usually related to underlying disease or complications.
